\section{Class \tkzClass{list\_point}}
\label{sec:class_list_point}

The variable \tkzVar{list\_point}{LP} holds a table used to store \tkzClass{list\_point} objects.
Its use is optional, and users are free to choose any variable name.

\medskip

However, for the sake of consistency and readability throughout the documentation,
it is recommended to use the predefined variable \tkzVar{list\_point}{LP}.
The function \tkzFct{tkz-elements}{init\_elements()} automatically reinitializes this table.

\medskip

The \tkzClass{list\_point} class represents an ordered list of points.
It is primarily intended as a lightweight container used to store and manipulate multiple points
produced by geometric constructions, intersections, or iterative algorithms.

Unlike geometric objects such as lines or circles, a \tkzClass{list\_point} object does not define intrinsic geometry.
Its purpose is to facilitate grouping, iteration, and data transfer between \tkzname{Lua} and \tkzname{TikZ}.

\medskip

A \tkzClass{list\_point} object behaves like a Lua table indexed by integers,
while providing a set of convenience methods adapted to geometric workflows.

\medskip

\subsection{Attributes}
\label{subsec:list_point_attributes}

\vspace{1em}
\bgroup
\small
\captionof{table}{\texttt{list\_point} attributes.}\label{list_point:attributes}
\begin{tabular}{lll}
\toprule
\textbf{Attribute} & \textbf{Meaning} & \texttt{Reference} \\
\midrule
\tkzAttr{list\_point}{n} &  [\ref{ssub:attr_list_point_n}] \\
\tkzAttr{list\_point}{items} & [\ref{ssub:attr_list_point_items}] \\
\bottomrule
\end{tabular}
\egroup


\subsubsection{Attribute \tkzAttr{list\_point}{n}}
\label{ssub:attr_list_point_n}

The attribute \tkzAttr{list\_point}{n} denotes the number of points stored in the list.
In practice, it corresponds to the Lua length operator \texttt{\#P} (array part).

\subsubsection{Attribute \tkzAttr{list\_point}{items}}
\label{ssub:attr_list_point_items}

The attribute \tkzAttr{list\_point}{items} refers to the ordered sequence of points stored in the list.
Individual points are accessed using standard Lua indexing, e.g. \texttt{P[i]}.


\subsection{Creating an object}
\label{subsec:list_point_create}

\subsubsection{Method \tkzMeth{list\_point}{new(...)}}
\label{ssub:meth_list_point_new}

Creates a new \tkzClass{list\_point} object.

In user code, the recommended constructor is the callable form \code{list\_point(...)}.
Internally, this calls the method \tkzMeth{list\_point}{new(...)}.

\medskip

A \tkzClass{list\_point} object can be created either as a local Lua variable
or as an entry of the table \tkzVar{list\_point}{LP}, which is reserved for storing
lists of points.

\begin{tkzexample}[code only]
local P = list_point()
LP.Q = list_point(z.A, z.B, z.C)
\end{tkzexample}

Both forms are valid and produce independent \tkzClass{list\_point} objects.

\medskip

Using a local variable is convenient for temporary computations or intermediate results.
However, for objects intended to be reused, transferred to \tkzname{TikZ},
or accessed outside a local scope, it is recommended to store them in the
\tkzVar{list\_point}{LP} table.

\medskip

The function \tkzFct{tkz-elements}{init\_elements()} automatically reinitializes
the \tkzVar{list\_point}{LP} table.


% --------------------------------------------------------------------
\subsection{Methods or Basic accessors}
\label{subsec:list_point_accessors}

\vspace{1em}
\bgroup
\small
\captionof{table}{\texttt{list\_point} methods.}\label{list_point:methods}
\begin{tabular}{ll}
\toprule
\texttt{Method}  & \texttt{Reference} \\
\midrule
\texttt{Constructor} &    \\
\midrule
\tkzMeth{list\_point}{new(...)}   & [\ref{ssub:meth_list_point_new}] \\
\midrule
\texttt{Methods Returning a Integer Number}&\\
\midrule
\tkzMeth{list\_point}{len()}      & [\ref{ssub:meth_list_point_len}] \\
\midrule
\texttt{Methods Returning a Point}&\\
\midrule\tkzMeth{list\_point}{get(i)}     & [\ref{ssub:meth_list_point_get}] \\
\tkzMeth{list\_point}{barycenter()}& [\ref{ssub:meth_list_point_barycenter}]\\
\midrule
\texttt{Methods Returning a List\_Point}&\\
\midrule
\tkzMeth{list\_point}{add(p)}     & [\ref{ssub:meth_list_point_add}] \\
\tkzMeth{list\_point}{extend(pl)} & [\ref{ssub:meth_list_point_extend}] \\
\tkzMeth{list\_point}{map(f)}     & [\ref{ssub:meth_list_point_map}] \\
\midrule
\tkzMeth{list\_point}{unpack()}   & [\ref{ssub:meth_list_point_unpack}] \\


\tkzMeth{list\_point}{clear()}    & [\ref{ssub:meth_list_point_clear}] \\
\tkzMeth{list\_point}{foreach(f)} & [\ref{ssub:meth_list_point_foreach}] \\

\tkzMeth{list\_point}{bbox()}     & [\ref{ssub:meth_list_point_bbox}] \\

\tkzMeth{list\_point}{to\_path()}& [\ref{ssub:meth_list_point_to_path}] \\
\bottomrule
\end{tabular}
\egroup

\subsubsection{Method \tkzMeth{list\_point}{len()}}
\label{ssub:meth_list_point_len}

Returns the number of points stored in the list.

\subsubsection{Method \tkzMeth{list\_point}{get(i)}}
\label{ssub:meth_list_point_get}

Returns the point at index \code{i}.

\subsubsection{Method \tkzMeth{list\_point}{unpack()}}
\label{ssub:meth_list_point_unpack}

Returns all points as separate values using \code{table.unpack}.
This method is useful when interfacing with functions expecting individual point arguments.

\begin{tkzexample}[latex=.5\textwidth]
\directlua{
init_elements()
LP = list_point()
LP:add(point(0,0))
LP:add(point(1,0))
LP:add(point(1,1))
z.A, z.B, z.C = LP:unpack()
}
\begin{center}
\begin{tikzpicture}
\tkzGetNodes
\tkzDrawPolygon(A,B,C)
\end{tikzpicture}
\end{center}
\end{tkzexample}

% --------------------------------------------------------------------
\subsection{Modification methods}
\label{subsec:list_point_mutators}

\subsubsection{Method \tkzMeth{list\_point}{add(p)}}
\label{ssub:meth_list_point_add}

Appends the point \code{p} to the list and returns the object itself.

\subsubsection{Method \tkzMeth{list\_point}{extend(pl)}}
\label{ssub:meth_list_point_extend}

Appends all points from another \tkzClass{list\_point} object \code{pl}.

\subsubsection{Method \tkzMeth{list\_point}{clear()}}
\label{ssub:meth_list_point_clear}

Removes all points from the list.


These methods modify the object in place and allow method chaining.

% --------------------------------------------------------------------
\subsection{Iteration helpers}
\label{subsec:list_point_iteration}

\subsubsection{Method \tkzMeth{list\_point}{foreach(f)}}
\label{ssub:meth_list_point_foreach}

Applies the function \code{f(p, i)} to each point \code{p} with its index \code{i}.

\subsubsection{Method \tkzMeth{list\_point}{map(f)}}
\label{ssub:meth_list_point_map}

Applies the function \code{f(p, i)} to each point and returns a new \tkzClass{list\_point} object
containing the results.


These methods support a functional programming style for geometric transformations.

% --------------------------------------------------------------------
\subsection{Geometric utilities}
\label{subsec:list_point_geometry}

\subsubsection{Method \tkzMeth{list\_point}{barycenter()}}
\label{ssub:meth_list_point_barycenter}

Returns the barycenter (centroid) of all points in the list.
If the list is empty, the method returns \code{nil}.

\subsubsection{Method \tkzMeth{list\_point}{bbox()}}
\label{ssub:meth_list_point_bbox}

Returns the bounding box of the point set as four numbers:
\[
  (x_{\min}, y_{\min}, x_{\max}, y_{\max})
\]


% --------------------------------------------------------------------
\subsection{TikZ output}
\label{subsec:list_point_output}

\subsubsection{Method \tkzMeth{list\_point}{to\_path()}}
\label{ssub:meth_list_point_to_path}

Converts the list of points into a TikZ path string of the form
\[
  (x_1,y_1) -- (x_2,y_2) -- \dots
\]
suitable for drawing with TikZ commands.


\medskip

This class provides a minimal yet effective bridge between \tkzname{Lua} computations and \tkzname{TikZ} drawing instructions.


\subsection{Examples}
\label{subsec:list_point_examples}

% --------------------------------------------------------------------
\subsubsection{Temporary list of points}
\label{ssub:list_point_example_local}

This example illustrates the use of a local \tkzClass{list\_point} object to store
temporary results inside a computation.

\begin{tkzexample}[code only]
local P = list_point()
for i = 1, 10 do
  P:add(z.A + i * (z.B - z.A) / 10)
end
local G = P:barycenter()
\end{tkzexample}

The list \code{P} is used only locally to compute the barycenter of a set of points.

% --------------------------------------------------------------------
\subsubsection{Storing construction results in \tkzVar{list\_point}{LP}}
\label{ssub:list_point_example_lp}

This example shows how to store a list of intersection points for later use.

\begin{tkzexample}[code only]
LP.I = list_point()
local X, Y = intersection(L.AB, C.circ)
LP.I:add(X)
LP.I:add(Y)
\end{tkzexample}

Storing the points in \tkzVar{list\_point}{LP} makes them accessible for drawing
or further constructions.

% --------------------------------------------------------------------
\subsubsection{Iterative construction}
\label{ssub:list_point_example_iterative}

A typical use case is the progressive construction of a family of points.

\begin{tkzexample}[code only]
LP.P = list_point()
for i = 0, 24 do
  local t = i / 24
  LP.P:add(L.AB:point(t))
end
\end{tkzexample}

The list \code{LP.P} contains a discrete sampling of a segment.

% --------------------------------------------------------------------
\subsubsection{Functional transformation}
\label{ssub:list_point_example_map}

This example uses the \tkzMeth{list\_point}{map} method to transform a set of points.

\begin{tkzexample}[code only]
local P = list_point(z.A, z.B, z.C)
LP.Q = P:map(function(p)
  return rotation_(p, tkz.tau / 3, z.O)
end)
\end{tkzexample}

The original list \code{P} is preserved, while \code{LP.Q} stores the transformed points.

% --------------------------------------------------------------------
\subsubsection{Conversion to a TikZ path}
\label{ssub:list_point_example_path}

This example illustrates the interaction between \tkzClass{list\_point} and TikZ.

\begin{tkzexample}[code only]
LP.P = list_point(z.A, z.B, z.C)
local path = LP.P:to_path()
\end{tkzexample}

The string returned by \tkzMeth{list\_point}{to\_path} can be passed directly
to TikZ drawing commands.


% --------------------------------------------------------------------
\subsubsection{From \tkzClass{list\_point} to \tkzClass{path} and drawing with \tkzNamePack{tkz-euclide}}
\label{ssub:list_point_example_points_path_draw}

This example shows a typical workflow: compute a list of points in \tkzname{Lua},
convert it into a \tkzClass{path}, and then draw the resulting points with
\tkzMeth{tkz-euclide}{tkzDrawPointsFromPath}.


\begin{tkzexample}[latex = .5\textwidth]
\directlua{
  init_elements()
  z.A = point(0, 0)
  z.B = point(8, 6)
  L.AB = line(z.A, z.B)
  % list of sampled points
  LP.P = list_point()
  for i = 0, 12 do
    LP.P:add(L.AB:point(i/12))
  end
  % convert list_point -> path
  PA.sample = path()
  for i = 1, LP.P:len() do
    PA.sample:add_point(LP.P:get(i))
  end
}
\begin{center}
\begin{tikzpicture}
\tkzGetNodes
% draw the points stored in the path
\tkzDrawPointsFromPath(PA.sample)
\end{tikzpicture}
\end{center}
\end{tkzexample}

The points are stored in \code{LP.P} for later reuse, while \code{PA.sample}
contains the corresponding TikZ path representation used for drawing.
